% ========================================================
% CORRECCIONES AL ENTREGABLE 1
% ========================================================
\section{Correcciones al Entregable 1}\label{sec:correcciones}

Las secciones anteriores reproducen el Entregable 1 tal como se entregó el 13 de septiembre
(tag \texttt{pc1} del repositorio). Esta sección recoge lo que cambió al pasar del diseño a la
implementación y lo que quedó pendiente de aquel informe. Los compromisos que se retiran se retiran
aquí, de forma expresa, y no en silencio.

\subsection{Estado de la asignación del caso}

El enunciado prevé que el docente asigne formalmente el caso de uso y el modelo. Al cierre de este
informe la asignación formal de \emph{NYC TLC + K-means} sigue sin confirmación escrita, y la
ausencia de respuesta no se toma como aprobación: el equipo continúa de forma provisional con la
propuesta de la Sección~\ref{sec:caso}, que es la que este entregable implementa y mide.

\subsection{Lo que cambió del diseño concurrente}

La Sección~\ref{sec:diseno} y el Algoritmo~\ref{alg:lloyd} describen el diseño de referencia de la
PC1. La implementación lo respeta en lo esencial (centroides de solo lectura durante la asignación,
acumuladores privados, barrera con \texttt{sync.WaitGroup}, reducción por sumas y conteos) y lo
precisa en cuatro puntos:

\begin{enumerate}
  \item \textbf{Acumuladores por bloque, no por worker.} El diseño de la PC1 daba a cada goroutine
        un acumulador $s_{wj}, n_{wj}$. En la implementación el acumulador pertenece al
        \emph{bloque} (\emph{chunk}) y la reducción suma los parciales en orden de índice de bloque.
        La razón es numérica: sumar en punto flotante no es asociativo, y si cada worker acumulara
        lo que le tocó, el resultado dependería del planificador. Con acumuladores por bloque, la
        versión concurrente produce exactamente el mismo resultado para cualquier cantidad de
        workers (Sección~\ref{sec:reproducibilidad}).
  \item \textbf{Reparto por canal, no en bloques contiguos por worker.} Los $P$ bloques fijos de la
        PC1 se reemplazan por muchos bloques pequeños que un canal reparte entre $P$ goroutines
        persistentes. El equilibrio de carga sale solo, y el tamaño de bloque pasa a ser la variable
        de experimento que la PC1 comprometió (Sección~\ref{sec:chunk}).
  \item \textbf{Inercia que se reporta.} El pseudocódigo de la PC1 devolvía los centroides
        actualizados junto con la inercia calculada con los anteriores. La implementación registra
        la inercia de cada iteración y, al final, la recalcula con los centroides finales en una
        pasada extra, para que el número publicado corresponda al modelo publicado.
  \item \textbf{Cluster vacío.} Conserva su centroide anterior; no se vuelve a sembrar, porque
        hacerlo con el punto más lejano haría que la corrida dependiera del orden de recorrido.
\end{enumerate}

\subsection{Compromisos que se retiran o se posponen}

\begin{itemize}
  \item \textbf{Variante con \texttt{sync.Mutex}.} La PC1 prometió una variante con acumuladores
        compartidos protegidos por mutex para cuantificar el costo de la contención. No se
        implementó. El contraste que sí se hizo es más informativo para la pregunta de fondo: un
        mutante del modelo en Promela con un acumulador compartido sin exclusión mutua, que Spin
        refuta con un contraejemplo (Sección~\ref{sec:promela}). Medir la contención de un mutex real
        queda para el Entregable 3.
  \item \textbf{Elección de $k$.} La PC1 propuso fijar $k$ comparando inercia y silueta sobre una
        muestra. Este entregable congela $k = \cifra{k}$ para todas las mediciones, sin ese estudio:
        el objeto de la PC2 es la concurrencia, y el speedup se mide con el mismo $k$ en ambas
        versiones. La elección de $k$ con criterio de validación interna queda para el Entregable 3.
  \item \textbf{Techo de Amdahl como predicción.} La PC1 estimó un techo de speedup suponiendo una
        fracción serial fija. Las mediciones muestran que la fracción serial efectiva crece con $P$
        (Sección~\ref{sec:karp}), así que ese techo no se publica como predicción: el límite no es
        una sección secuencial fija sino el costo de coordinar.
\end{itemize}

\subsection{Compromisos que se cumplen}

Se midió el escalamiento fuerte y el débil (Secciones~\ref{sec:fuerte} y~\ref{sec:debil}), el
tamaño de bloque se trató como variable del experimento (Sección~\ref{sec:chunk}), la equivalencia
entre versiones se comprobó con tolerancia explícita y, además, comparando las asignaciones
(Sección~\ref{sec:equivalencia}), la lectura del CSV se midió aparte de las iteraciones, y la
barrera se modeló en Promela y se verificó con Spin.

\subsection{Revisión de las features}

Las features de la Tabla~\ref{tab:features} se revisaron al implementar el lector en Go. Las cuatro
componentes seno y coseno no se estandarizan por separado, porque eso deformaría el círculo y
rompería la vecindad entre las 23:00 y la 01:00; las dos transformaciones \texttt{log1p} sí llevan
z-score. Las varianzas medidas son 0,389 y 0,487 para hora, 0,509 y 0,487 para día, y 1,0 para
duración y distancia, así que los cuatro conceptos pesan parecido en la distancia euclídea. Duración y
distancia están correlacionadas ($r = 0{,}833$) y se conservan las dos a propósito: dos viajes de la
misma distancia y distinta duración son tipos de viaje distintos, que es justo lo que el caso quiere
distinguir.

\subsection{Participación del tercer integrante}

El informe de la PC1 registró que el tercer integrante no tenía commits en el repositorio. Desde
entonces tiene acceso como colaborador (usuario \texttt{ElJulioGG}). Su aporte a este entregable se
documenta en la Sección~\ref{sec:github}.
